\begin{flushright}
{\large ПРИЛОЖЕНИЕ 1}\\[0.5em]
Бланк индивидуального задания на практику студента
\end{flushright}

\vspace{1em}

\begin{center}
\textbf{\MakeUppercase{\UniversityName}}
\end{center}

\noindent\textbf{Направление:} \DirectionName\\
\textbf{ОП:} \ProgramName

\vspace{1.5em}

\begin{center}
{\Large\bfseries ИНДИВИДУАЛЬНОЕ ЗАДАНИЕ}\\[0.3em]
на научно-исследовательскую практику
\end{center}

\vspace{1em}

\noindent\textbf{Студент:} \StudentName \hfill \textbf{номер группы:} \StudentGroup

\noindent\textbf{Тема задания:} \PracticeTheme

\noindent\textbf{Сроки прохождения практики:} \PracticeStart{} -- \PracticeEnd

\noindent\textbf{Место прохождения практики:} \PracticePlace

\vspace{0.75em}

\noindent\textbf{1. Виды работ и требования к их выполнению:}
\begin{enumerate}[leftmargin=1.5em]
\item Провести анализ предметной области и сформулировать задачу как прототип системы поддержки принятия решений для рекомендации следующей культуры.
\item Подготовить и проверить воспроизводимый пайплайн: этапы очистки данных, построения оконного датасета и разделения на обучающую, валидационную и тестовую выборки.
\item Исследовать baseline-подходы и выбрать финальную модель для предсказания.
\item Оценить качество модели и промежуточных экспериментов, включая вариант с дополнительными признаками и эксперимент \texttt{CatBoost + transition graph}.
\item Подготовить краткий отчет по практике и комплект приложений с опорой на фактические результаты дипломной работы.
\end{enumerate}

\noindent\textbf{2. Виды отчетных материалов и требования к их оформлению:}
\begin{enumerate}[leftmargin=1.5em]
\item Отчет по практике с описанием цели, постановки задачи, данных, этапов подготовки выборки, модели предсказания и промежуточных экспериментов.
\item Дневник практики с фиксацией этапов работ по периодам прохождения практики.
\item Результаты экспериментов, таблицы и графики.
\end{enumerate}

\vspace{1.5em}

\begin{center}
\textbf{3. ПЛАН-ГРАФИК}
\end{center}

\small
\setlength{\tabcolsep}{3pt}
\begin{center}
\begin{tabularx}{\textwidth}{|C{0.8cm}|Y|C{2.1cm}|Y|Y|}
\textbf{№ этапа} & \textbf{Наименование этапа} & \textbf{Срок завершения} & \textbf{Виды работ} & \textbf{Форма отчетности} \\
\hline
1 & Постановка задачи и обзор материалов & 20.02.2026 & Анализ темы диплома, структуры работы и постановки задачи рекомендации & Рабочие заметки, уточненная формулировка темы \\
\hline
2 & Подготовка данных и формирование выборки & 27.02.2026 & Очистка данных, оконное представление, таблица признаков & Описание датасета и подготовленной выборки \\
\hline
3 & Baseline-подходы и разделение датасета & 10.03.2026 & Анализ простых baseline-моделей, сплит датасета по \texttt{CSBID}, контроль воспроизводимости & Табличное сравнение baseline-подходов \\
\hline
4 & Финальная baseline-модель & 23.03.2026 & Проверка CatBoost, состава признаков и ключевых метрик качества & Сводка качества модели, описание признаков \\
\hline
5 & Промежуточные эксперименты & 03.04.2026 & Анализ ветки с дополнительными признаками и эксперимента \texttt{CatBoost + transition graph} & Сравнительные таблицы и вывод о выборе финальной модели \\
\hline
6 & Оформление итоговых материалов практики & 14.04.2026 & Подготовка краткого отчета, дневника & Комплект приложений и итоговый отчет \\
\hline
\end{tabularx}
\end{center}
\setlength{\tabcolsep}{6pt}
\normalsize

\vspace{1.5em}

\noindent Задание утверждено на заседании кафедры \underline{\hspace{7cm}}

\noindent (протокол от «\underline{\hspace{0.8cm}}» \underline{\hspace{3cm}} 20\underline{\hspace{0.5cm}} г. №\underline{\hspace{1cm}}).

\vspace{1em}

\noindent Руководитель практики \hfill \underline{\hspace{3cm}} \hfill \SupervisorName

\vspace{1.5em}

\noindent Задание принял к исполнению «11» февраля 2026 г.

\vspace{1.5em}

\noindent\hfill \underline{\hspace{4cm}} \hfill \StudentName
